\section{Phase decomposition and exact pattern counts}\label{sec:phase}

Fix disjoint sets $A$ and $B$ with $|A|=m$ and $|B|=n$.  The graph
$K_{m,n}$ has vertex set $A\sqcup B$ and every edge joins $A$ to $B$.  Write
\[
 X_{m,n}^{(d)}=\{x\in(A\sqcup B)^{\Z^d}:x_v\sim x_{v+e_i}
 \text{ for all }v\in\Z^d,\ 1\leq i\leq d\}.
\]
Let
\[
 \chi(v)=v_1+\cdots+v_d\pmod 2,
 \qquad \E=\ker\chi,
 \qquad \OO=\Z^d\setminus\E.
\]

\begin{lemma}[Intrinsic phase]\label{lem:phase}
Every $x\in X_{m,n}^{(d)}$ has exactly one of the following orientations:
\[
x_{\E}\subset A,\quad x_{\OO}\subset B,
\qquad\text{or}\qquad
x_{\E}\subset B,\quad x_{\OO}\subset A.
\]
If $\omega(x)=1$ in the first case and $-1$ in the second, then
$\omega(\sigma^v x)=(-1)^{\chi(v)}\omega(x)$.
\end{lemma}

\begin{proof}
Choose the part containing $x_0$.  Along every nearest-neighbour edge the
part changes.  The parity of the length of a lattice path from $0$ to $v$ is
$\chi(v)$, so the part at $v$ is forced and is independent of the chosen
path.  Translation by an odd vector exchanges the two parity classes, while
translation by an even vector preserves them.
\end{proof}

\begin{proposition}[All finite shapes]\label{prop:patterns}
The number $N_{m,n}(F)$ of patterns on a finite set $F$ that extend to
$X_{m,n}^{(d)}$ satisfies $N_{m,n}(\varnothing)=1$.  For nonempty $F$ it is
\[
 N_{m,n}(F)=
 m^{|F\cap\E|}n^{|F\cap\OO|}
 +n^{|F\cap\E|}m^{|F\cap\OO|}.
 \tag{2.1}\label{eq:finite-count}
\]
In particular,
\[
 h_{\rm top}(X_{m,n}^{(d)})=\frac12\log(mn).
 \tag{2.2}\label{eq:entropy}
\]
\end{proposition}

\begin{proof}
The empty restriction is unique.  If $F$ is nonempty, the two alternatives
in \cref{lem:phase} are disjoint on $F$.  In the first phase the sites in
$F\cap\E$ have $m$ choices and those in $F\cap\OO$ have $n$ choices; the
second phase reverses the roles.  Every such restriction extends: choose
arbitrary symbols from the phase-prescribed target part at every site outside
$F$.  Completeness of the bipartite target makes the resulting global
configuration valid.  Summing the two phase counts proves
\eqref{eq:finite-count}.

Take rectangular Følner boxes.  They are connected and their two parity
classes have sizes $|F|/2+O(1)$.  Applying \eqref{eq:finite-count}, taking
$|F|^{-1}\log$, and passing to the limit gives \eqref{eq:entropy}.
\end{proof}

The global choice of phase contributes only a bounded factor, while each
even--odd pair contributes $mn$ choices.  Importantly, disconnected pieces
of $F$ cannot choose phases independently: global extendibility couples them
through \cref{lem:phase}.  The conjugacy in the next section makes that
pairing local.

Let $X^+$ be the phase component with $A$ on $\E$, and let $X^-$ be the
other component.  Although neither component is invariant under the full
$\Z^d$ action, both are invariant under the subgroup $\E$.

\begin{proposition}[Full-shift model for a phase]\label{prop:phase-full}
As an $\E$-system, $X^+$ is conjugate to the full shift
$(A\times B)^{\E}$.  The map is
\[
 \Theta_+(x)_v=(x_v,x_{v+e_1}),\qquad v\in\E.
 \tag{2.3}\label{eq:phase-pair}
\]
The component $X^-$ has the same full-shift model after translating the
anchors by $e_1$.
\end{proposition}

\begin{proof}
The edges $(v,v+e_1)$ with $v\in\E$ are disjoint and cover $\Z^d$.
Consequently \eqref{eq:phase-pair} records every input coordinate exactly
once.  Conversely, an arbitrary labelling of $\E$ by $A\times B$ can be
unpacked independently on these dimers and gives a valid point of $X^+$,
because every lattice edge joins an $A$ symbol to a $B$ symbol.  Packing and
unpacking are continuous and commute with translations in $\E$.  For $X^-$,
the $A$-anchors form the coset $e_1+\E$; translating this coset back to
$\E$ gives the identical construction.
\end{proof}

An odd translation exchanges $X^+$ and $X^-$.  It follows that every
full-action invariant probability is the equal mixture of an
$\E$-invariant probability on one component and its odd translate.  This
observation will turn the equilibrium problem into the ordinary full-shift
variational problem on the dimer alphabet.
